% ================================================================
%  Chapter 4 — High Level Design (KANDA)
%  Knowledge-driven Autonomous Navigation and Decision-making Agent
% ================================================================
\chapter{High Level Design}
\label{ch:hld}

This chapter explains how we designed KANDA to solve the problems from Chapter 1 while meeting the requirements from Chapter 3. Rather than start with formal specifications, we describe the real design choices and their rationale. The ESP32 handles time-critical sensing and motor control. The Raspberry Pi handles slower language-model tasks. Safety comes from three independent validation layers that check every command before motors run. We learned early that language models produce garbage sometimes, and you don't want that reaching the motors.

%-------------------------------------------------------------------
\section{Design Considerations}
\label{sec:design}

%-------------------------------------------------------------------
\subsection{General Considerations}
\label{subsec:general}

The following considerations shaped the high-level design of KANDA:

\textbf{Safety before motion:} Language models sometimes produce invalid output. Malformed JSON, made-up action names, speed values outside the valid range. Raw LLM output cannot go directly to motors. Every command must pass strict checks. Is it valid JSON? Is the action in the whitelist? Is the speed between 0 and 255? If any check fails, the robot stops.

\textbf{Layered architecture:} All components organize into independent layers. The bottom layer handles raw sensing and motor control (ESP32). The middle layer manages message passing and validation (Raspberry Pi). The top layer does reasoning (Groq, NVIDIA NIM). Each layer has a narrow interface: telemetry lines and JSON lines only. This makes debugging easier because you know where to look when something breaks.

\textbf{Timing mismatch:} The ESP32 samples sensors every 100 milliseconds. Cloud API calls take 1.5 to 4 seconds. Waiting for every cloud response would freeze the robot. Instead, the ESP32 runs local obstacle-avoidance logic independently. When the Pi goes silent (network lag or API timeout), the ESP32 avoids obstacles on its own. Not intelligent, but safe.

\textbf{Structured commands:} Instead of sending natural language to the firmware ("move forward carefully"), we send only JSON with two fields: \texttt{action} (must be one of seven words) and \texttt{speed} (must be 0--255). This way the firmware parser is dead simple and cannot be confused.

\textbf{Observability:} The OLED screen shows the robot's current state (listening, thinking, acting, searching). You can watch the robot without opening a terminal. The Pi logs all events as well. This combination saved countless debugging hours.

\textbf{Graceful degradation:} Loss of Wi-Fi, API downtime, or Pi crashes don't freeze the robot. The firmware detects silence from the Pi and switches to obstacle-avoidance mode. Not smart, but the robot stays alive.

%-------------------------------------------------------------------
\subsection{Development Methods}
\label{subsec:dev_methods}

KANDA was built incrementally: each phase works on its own before we added the next layer.

\begin{description}[nosep]
  \item[Phase 1, Foundations:] Read papers, sketch requirements, design the basic architecture (Chapters 1-3). Nothing physical yet.
  
  \item[Phase 2, Embodiment layer:] Get the ESP32 working: read ultrasonic sensors, drive motors via the H-bridge, display status on the OLED. By itself, it could avoid obstacles. No Pi, no AI, just firmware. This baseline took longer than expected because voltage dividers and sensor crosstalk are tricky.
  
  \item[Phase 3, Intelligence layer:] Add the Raspberry Pi running Python, connect to Groq (text) and NVIDIA NIM (vision), build prompts, send JSON commands to the ESP32. The firmware gets an AI\_MODE for good commands and an AUTO\_MODE fallback when the Pi is silent. The 3-second timeout proved critical for safety.
  
  \item[Phase 4, Multimodal agent:] Add camera, microphone, wake-word detection, text-to-speech, Telegram bot, and a 7-state event machine. The OLED gets animated faces. The Pi understands images and voice. This phase focuses on integration. ReAct search with semantic memory prevents revisiting areas.
\end{description}

%-------------------------------------------------------------------
\section{Architectural Strategies}
\label{sec:arch_strat}

%-------------------------------------------------------------------
\subsection{Programming Language and Runtime Selection}
\label{subsec:lang}

\textbf{Arduino-compatible C++} targets the ESP32 firmware: direct register and peripheral access, deterministic sampling intervals, mature motor-control patterns, and small binary size~\cite{espressif2023esp32}.

\textbf{Python~3} targets the Raspberry Pi intelligence layer: \texttt{pyserial} for UART communication, and fast iteration on reasoning. The system uses two cloud APIs: Groq Llama 3.3 70B for text reasoning and motion planning, and NVIDIA NIM Llama 3.2 11B Vision for visual scene understanding. We use \texttt{openWakeWord} (offline, no API) for wake detection, Google Text-to-Speech (free, no key) for voice output, and \texttt{picamera2} for image capture.

\textbf{JSON line protocol} over UART at 115,200 baud provides human-readable debugging, explicit message boundaries (newline-terminated lines), and straightforward parsing on both ends without custom framing beyond newline delimiters.

%-------------------------------------------------------------------
\subsection{Future Extensions}
\label{subsec:future}

Planned extensions beyond the baseline include on-device telemetry buffering, SLAM integration, depth sensors, and local-model experiments (Gemma 2B). All future work maintains the validation layer before actuation~\cite{huang2023grounded}.

%-------------------------------------------------------------------
\subsection{Safety Validation Pipeline}
\label{subsec:safety_pipeline}

A central design principle is \textbf{defence-in-depth}: no single layer alone protects the physical plant. Figure~\ref{fig:safety_pipeline} illustrates the three-layer safety pipeline that intercepts every command between the LLM and the motors.

\begin{figure}[H]
\centering
\begin{tikzpicture}[
  block/.style={draw, thick, rounded corners=3pt, minimum height=0.9cm,
                minimum width=3.2cm, align=center, font=\footnotesize},
  arr/.style={-{Stealth[length=2.5mm]}, thick},
  node distance=0.6cm
]
\node[block, fill=gray!10] (prompt) {Specification\\(Body-Context Prompt)};
\node[block, fill=yellow!12, below=of prompt] (llmout) {LLM Output\\(JSON candidate)};
\node[block, fill=orange!15, below=of llmout] (validator) {Runtime Validator\\(whitelist + clamp)};
\node[block, fill=green!12, below=of validator] (uart) {UART Transmission\\(validated JSON)};
\node[block, fill=red!10, below=of uart] (firmware) {Firmware Safety\\(obstacle stop)};
\node[block, fill=blue!8, below=of firmware] (motors) {Motor Actuation};

\draw[arr] (prompt) -- node[right, font=\tiny]{action vocab, safety rules} (llmout);
\draw[arr] (llmout) -- node[right, font=\tiny]{may be malformed} (validator);
\draw[arr] (validator) -- node[right, font=\tiny]{always valid dict} (uart);
\draw[arr] (uart) -- (firmware);
\draw[arr] (firmware) -- node[right, font=\tiny]{stop if front $<$15\,cm} (motors);

% Rejection arrows
\node[right=2.5cm of validator, font=\tiny\itshape, text=red!70!black] (reject) {violation $\to$ stop};
\draw[-{Stealth[length=2mm]}, dashed, red!60!black] (validator.east) -- (reject.west);
\end{tikzpicture}
\caption{Three-layer safety validation pipeline: specification-level grounding in the prompt, runtime enforcement on the Pi, and device-level emergency stop in ESP32 firmware.}
\label{fig:safety_pipeline}
\end{figure}

%-------------------------------------------------------------------
\section{System Architecture}
\label{sec:sysarch}

KANDA is structured as a \textbf{three-tier logical architecture}: intelligence on the Raspberry Pi, a duplex UART link, and embodiment on the ESP32 interfacing to sensors, drivers, and display. Figure~\ref{fig:sysarch} summarises the layering.

\begin{figure}[htbp!]
\centering
\resizebox{0.92\textwidth}{!}{%
\begin{tikzpicture}[
  layer/.style={draw, thick, rounded corners=4pt, minimum width=11.8cm, minimum height=1.35cm, align=center, font=\small},
  halfbox/.style={draw, thick, rounded corners=4pt, minimum width=5.6cm, minimum height=1.05cm, align=center, font=\footnotesize},
  arr/.style={-{Stealth[length=2.5mm]}, thick},
  lbl/.style={font=\scriptsize, fill=white, inner sep=2pt},
  node distance=0.65cm
]

% Layer 1: Intelligence
\node[layer, fill=blue!8] (intel) {
  \textbf{Intelligence layer (Raspberry Pi)}\\[3pt]
  {\footnotesize Body context \quad$\bullet$\quad Groq/NVIDIA clients \quad$\bullet$\quad Safety validator \quad$\bullet$\quad UART bridge}
};

% UART link (narrow band)
\node[layer, fill=yellow!15, minimum height=0.85cm, font=\footnotesize] (uart) [below=of intel] {
  \textbf{UART serial}\quad 115,200 baud\quad newline-delimited JSON \& telemetry
};

% Embodiment: two halves
\node[halfbox, fill=orange!10, below left=0.65cm and 0.25cm of uart] (embL) {
  \textbf{Control \& parsing}\\[2pt]
  {\scriptsize JSON $\rightarrow$ PWM \& direction}
};
\node[halfbox, fill=orange!10, below right=0.65cm and 0.25cm of uart] (embR) {
  \textbf{Sensing \& display}\\[2pt]
  {\scriptsize HC-SR04 \quad OLED status}
};

\node[layer, fill=orange!6, minimum width=11.8cm, minimum height=0.95cm, font=\footnotesize] (embbar) at ($(embL.south)!0.5!(embR.south) + (0,-0.65)$) {
  \textbf{Embodiment layer (ESP32)} --- dual-mode: AI\_MODE when valid JSON arrives; AUTO\_MODE after watchdog timeout
};

% Physical
\node[layer, fill=gray!12, below=0.55cm of embbar] (phys) {
  \textbf{Physical plant}\\[2pt]
  {\footnotesize HC-SR04 sensors \quad$\bullet$\quad TB6612FNG \quad$\bullet$\quad DC motors \quad$\bullet$\quad power \& mechanics}
};

\draw[arr] (intel) -- node[lbl, right] {Validated JSON downlink} (uart);
\draw[arr] (uart) -- node[lbl, left, pos=0.45] {Telemetry uplink} (intel);

\draw[arr] (uart.south) -- ++(0,-0.12) -| (embL.north);
\draw[arr] (uart.south) -- ++(0,-0.12) -| (embR.north);

\draw[arr] (embL.south) -- ++(0,-0.08) -| (embbar.155);
\draw[arr] (embR.south) -- ++(0,-0.08) -| (embbar.25);

\draw[arr] (embbar) -- node[lbl, right] {PWM \& sensing} (phys);

\end{tikzpicture}%
}% end resizebox
\caption{High-level layered architecture of KANDA}
\label{fig:sysarch}
\end{figure}

Fig.~\ref{fig:sysarch} shows how responsibilities are split. The \textit{intelligence layer} performs hardware-aware prompting, calls Groq (text) and NVIDIA NIM (vision), validates JSON commands, and transmits safe motor intents over UART. The \textit{UART link} carries newline-terminated ASCII lines in both directions. The \textit{embodiment layer} parses commands, drives motors, samples ultrasonic sensors, refreshes the OLED, and enforces mode logic including fallback to AUTO\_MODE. The \textit{physical plant} captures sensors, actuators, and power paths documented in hardware chapters.

%-------------------------------------------------------------------
\section{Data Flow Diagrams}
\label{sec:dfd}

%-------------------------------------------------------------------
\subsection{Data Flow Diagram Level~0, Overall System}
\label{subsec:dfd0}

The Level~0 DFD places KANDA at the centre of supervision, cloud inference, and the physical deployment environment.

\begin{figure}[htbp!]
\centering
\resizebox{0.88\textwidth}{!}{%
\begin{tikzpicture}[
  entity/.style={draw, thick, rectangle, rounded corners=3pt, minimum width=3.2cm, minimum height=1cm, align=center, font=\small},
  process/.style={draw, thick, circle, minimum size=2.9cm, align=center, font=\small},
  arr/.style={-{Stealth[length=2.5mm]}, thick},
  lbl/.style={font=\scriptsize, fill=white, inner sep=2pt},
  node distance=1.4cm
]

\node[entity, fill=blue!10] (op) {\textbf{Operator}\\{\footnotesize (supervision)}};

\node[process, fill=green!12, below=2cm of op] (sys) {\textbf{KANDA}\\{\footnotesize embodied agent}};

\node[entity, fill=purple!10, below left=1.9cm and 2.5cm of sys] (gem) {\textbf{Groq + NVIDIA NIM}\\{\footnotesize (HTTPS)}};

\node[entity, fill=gray!15, below right=1.9cm and 2.5cm of sys] (env) {\textbf{Environment}\\{\footnotesize (workspace)}};

\draw[arr] (op) -- node[lbl, left] {start/stop, config} (sys);
\draw[arr] (sys) -- node[lbl, right] {status, logs} (op);

\draw[arr] ([xshift=-0.35cm]sys.west) -- node[lbl, above, sloped] {prompt / completion} ([yshift=0.12cm]gem.east);
\draw[arr] ([yshift=-0.12cm]gem.east) -- node[lbl, below, sloped] {JSON text} ([xshift=-0.35cm]sys.west);

\draw[arr] ([xshift=0.35cm]sys.east) -- node[lbl, above, sloped] {motion} ([yshift=0.12cm]env.west);
\draw[arr] ([yshift=-0.12cm]env.west) -- node[lbl, below, sloped] {layout / obstacles} ([xshift=0.35cm]sys.east);

\end{tikzpicture}%
}% end resizebox
\caption{Data flow diagram Level~0, overall system}
\label{fig:dfd0}
\end{figure}

Fig.~\ref{fig:dfd0} abstracts KANDA as a single process. The \textit{operator} starts or stops runs and may adjust configuration while receiving status and logs. The system exchanges HTTPS messages with \textit{Groq and NVIDIA NIM}: prompts ascend and structured completions descend. Interaction with the \textit{environment} captures that motion changes the scene layout---including obstacle distances reflected indirectly through onboard sensing inside the system boundary.

%-------------------------------------------------------------------
\subsection{Data Flow Diagram Level~1, Embodiment and Intelligence}
\label{subsec:dfd1}

The Level~1 DFD decomposes KANDA into two cooperating sub-systems linked by the UART channel.

\begin{figure}[htbp!]
\centering
\resizebox{0.94\textwidth}{!}{%
\begin{tikzpicture}[
  entity/.style={draw, thick, rectangle, rounded corners=3pt, minimum width=3cm, minimum height=0.75cm, align=center, font=\small, fill=purple!10},
  process/.style={draw, thick, rectangle, rounded corners=6pt, minimum width=5.4cm, minimum height=0.52cm, align=center, font=\small},
  chan/.style={draw, thick, rectangle, rounded corners=2pt, minimum width=2.2cm, minimum height=4.6cm, align=center, font=\footnotesize, fill=yellow!14},
  arr/.style={-{Stealth[length=2mm]}, thick},
  lbl/.style={font=\scriptsize, fill=white, inner sep=2pt},
  node distance=0.32cm
]

\node[entity] (gem) {\textbf{Groq + NVIDIA NIM}};

% UART channel (centre spine)
\node[chan, below=1cm of gem] (uart) {\textbf{UART}\\[4pt]{\scriptsize duplex}};

% LEFT: Embodiment column (ESP32)
\node[process, fill=orange!14, font=\small\bfseries, left=2.15cm of uart, anchor=east, minimum width=5.5cm] (e0) {Embodiment (ESP32)};
\node[process, fill=orange!6, below=0.2cm of e0, font=\scriptsize, minimum width=5.5cm] (e1) {1.1 Sample ultrasonic ranges};
\node[process, fill=orange!6, below=0.12cm of e1, font=\scriptsize, minimum width=5.5cm] (e2) {1.2 Format telemetry line};
\node[process, fill=orange!6, below=0.12cm of e2, font=\scriptsize, minimum width=5.5cm] (e3) {1.3 Parse JSON command};
\node[process, fill=orange!6, below=0.12cm of e3, font=\scriptsize, minimum width=5.5cm] (e4) {1.4 PWM, direction, OLED};

% RIGHT: Intelligence column (Pi)
\node[process, fill=blue!12, font=\small\bfseries, right=2.15cm of uart, anchor=west, minimum width=5.5cm] (i0) {Intelligence (Raspberry Pi)};
\node[process, fill=blue!6, below=0.2cm of i0, font=\scriptsize, minimum width=5.5cm] (i1) {2.1 Read telemetry};
\node[process, fill=blue!6, below=0.12cm of i1, font=\scriptsize, minimum width=5.5cm] (i2) {2.2 Build hardware-aware prompt};
\node[process, fill=blue!6, below=0.12cm of i2, font=\scriptsize, minimum width=5.5cm] (i3) {2.3 Invoke Groq/NIM};
\node[process, fill=blue!6, below=0.12cm of i3, font=\scriptsize, minimum width=5.5cm] (i4) {2.4 Validate JSON command};
\node[process, fill=blue!6, below=0.12cm of i4, font=\scriptsize, minimum width=5.5cm] (i5) {2.5 Transmit validated command};

% Cloud APIs to intelligence (HTTPS inference)
\draw[arr] (gem.south) -- +(0,-0.35) -| node[lbl, pos=0.28, above] {HTTPS} (i3.north);

% Telemetry uplink: embodiment -> uart -> intelligence
\draw[arr] (e2.east) -- node[lbl, above, pos=0.4] {telemetry} (uart.west);

\draw[arr] (uart.east) -- node[lbl, above, pos=0.4] {telemetry} (i1.west);

% Command downlink: intelligence -> uart -> embodiment
\draw[arr] (i5.west) -- node[lbl, below, pos=0.4] {JSON} (uart.east);

\draw[arr] (uart.west) -- node[lbl, below, pos=0.4] {JSON} (e3.east);

\end{tikzpicture}%
}% end resizebox
\caption{Data flow diagram Level~1, embodiment and intelligence sub-systems}
\label{fig:dfd1}
\end{figure}

Fig.~\ref{fig:dfd1} separates real-time embodiment (left) from asynchronous intelligence (right). Telemetry formatted on the ESP32 flows upward to the Raspberry Pi; validated JSON commands flow downward after prompt construction, cloud inference, and validation. HTTPS exchanges with Groq and NVIDIA NIM attach to the intelligence column (expanded at Level~2).

%-------------------------------------------------------------------
\subsection{Data Flow Diagram Level~2, Intelligence Pipeline}
\label{subsec:dfd2_intel}

\begin{figure}[htbp!]
\centering
\resizebox{0.82\textwidth}{!}{%
\begin{tikzpicture}[
  proc/.style={draw, thick, rectangle, rounded corners=4pt, minimum width=8.8cm, minimum height=0.78cm, align=center, font=\footnotesize, fill=blue!8},
  ext/.style={draw, thick, rectangle, rounded corners=3pt, minimum width=3.4cm, minimum height=0.55cm, align=center, font=\footnotesize, fill=purple!10},
  arr/.style={-{Stealth[length=2mm]}, thick},
  lbl/.style={font=\scriptsize, fill=white, inner sep=1pt},
  node distance=0.38cm
]

\node[ext] (uartin) {\textbf{UART} {\scriptsize (telemetry in)}};

\node[proc, below=of uartin] (p1) {\textbf{2.1 Read telemetry} (parse distances and state from ASCII line)};

\node[proc, below=of p1] (p2) {\textbf{2.2 Build prompt} (inject ranges, legal actions, speed bounds, rules)};

\node[proc, below=of p2] (p3) {\textbf{2.3 Invoke Groq/NVIDIA} (low temperature, bounded output tokens)};

\node[ext, below=of p3] (gem) {\textbf{Groq / NVIDIA NIM APIs}};

\node[proc, below=of gem] (p4) {\textbf{2.4 Validate} (schema, whitelist, clamp speed, safe stop on fault)};

\node[ext, below=of p4] (uartout) {\textbf{UART} {\scriptsize (JSON command out)}};

\draw[arr] (uartin) -- node[lbl, right] {line} (p1);
\draw[arr] (p1) -- node[lbl, right] {structs} (p2);
\draw[arr] (p2) -- node[lbl, right] {prompt} (p3);
\draw[arr] (p3) -- node[lbl, right] {HTTPS} (gem);
\draw[arr] (gem) -- node[lbl, right] {raw JSON} (p4);
\draw[arr] (p4) -- node[lbl, right] {safe command} (uartout);

\end{tikzpicture}%
}% end resizebox
\caption{Data flow diagram Level~2, intelligence pipeline on the Raspberry Pi}
\label{fig:dfd2a}
\end{figure}

Fig.~\ref{fig:dfd2a} details the Pi-side pipeline implemented by \path|vision_module/main.py| and companion modules: telemetry ingestion, hardware-aware prompt construction, Groq/NVIDIA invocation, validation, and transmission of a single safe JSON line per cycle.

%-------------------------------------------------------------------
\subsection{Data Flow Diagram Level~2, Embodiment Pipeline}
\label{subsec:dfd2_emb}

\begin{figure}[htbp!]
\centering
\resizebox{0.82\textwidth}{!}{%
\begin{tikzpicture}[
  proc/.style={draw, thick, rectangle, rounded corners=4pt, minimum width=8.8cm, minimum height=0.78cm, align=center, font=\footnotesize, fill=orange!8},
  ext/.style={draw, thick, rectangle, rounded corners=3pt, minimum width=3.4cm, minimum height=0.55cm, align=center, font=\footnotesize, fill=gray!14},
  arr/.style={-{Stealth[length=2mm]}, thick},
  lbl/.style={font=\scriptsize, fill=white, inner sep=1pt},
  node distance=0.38cm
]

\node[ext] (sns) {\textbf{Sensors} {\scriptsize (HC-SR04)}};

\node[proc, below=of sns] (p1) {\textbf{1.1 Trigger and capture echoes} (compute distances in cm)};

\node[proc, below=of p1] (p2) {\textbf{1.2 Compose telemetry} (ASCII line to companion)};

\node[ext, below=of p2] (uartup) {\textbf{UART TX} {\scriptsize (to Pi)}};

\node[ext, below=0.45cm of uartup] (uartdn) {\textbf{UART RX} {\scriptsize (from Pi)}};

\node[proc, below=of uartdn] (p3) {\textbf{1.3 Parse JSON command} (ArduinoJson parser)};

\node[proc, below=of p3] (p4) {\textbf{1.4 Map action to motor primitives} (PWM and direction pins)};

\node[ext, below=of p4] (mot) {\textbf{Motors \& OLED}};

\draw[arr] (sns) -- node[lbl, right] {echo timing} (p1);
\draw[arr] (p1) -- node[lbl, right] {ranges} (p2);
\draw[arr] (p2) -- node[lbl, right] {telemetry} (uartup);

\draw[arr] (uartdn) -- node[lbl, right] {JSON line} (p3);
\draw[arr] (p3) -- node[lbl, right] {intent} (p4);
\draw[arr] (p4) -- node[lbl, right] {drive / text} (mot);

\end{tikzpicture}%
}% end resizebox
\caption{Data flow diagram Level~2, embodiment pipeline on the ESP32}
\label{fig:dfd2b}
\end{figure}

Fig.~\ref{fig:dfd2b} expands embodiment-side processing: ranging, telemetry uplink, inbound JSON parsing, mapping discrete actions to TB6612FNG outputs, and operator-visible OLED updates. Watchdog and AUTO\_MODE transitions are policy logic around these flows.

%-------------------------------------------------------------------
\section{Summary}
\label{sec:ch4summary}

This chapter presented the high-level design of KANDA. Design considerations emphasised validation before actuation~\cite{huang2023grounded}, affordance-grounded command formats~\cite{ahn2022saycan}, layered embodied structure~\cite{brooks1986robust}, and resilience to stochastic cloud latency~\cite{vemprala2024chatgpt}. Development followed phased delivery from embodiment to intelligence, with a documented path toward multimodal extensions~\cite{zeng2023socratic,wu2023tidybot}.

Architectural strategies justified C++ on the ESP32 and Python on the Raspberry Pi, linked by a newline-delimited JSON and telemetry UART protocol~\cite{espressif2023esp32}. Figure~\ref{fig:sysarch} summarised the intelligence, UART, embodiment, and physical tiers. Data flow diagrams at Level~0 (\ref{fig:dfd0}), Level~1 (\ref{fig:dfd1}), and Level~2 (\ref{fig:dfd2a}, \ref{fig:dfd2b}) traced information from sensors and cloud inference to motor outputs. Chapter~5 refines these sub-systems into detailed module-level design aligned with the implementation directories.
